\documentclass[11pt,letterpaper]{article}

% ── Geometry ──────────────────────────────────────────────────────────────────
\usepackage[margin=1in]{geometry}

% ── Pagination ────────────────────────────────────────────────────────────────
\usepackage[all]{nowidow}
\brokenpenalty=10000
\raggedbottom

% ── Fonts ─────────────────────────────────────────────────────────────────────
\usepackage[T1]{fontenc}
\usepackage{newpxtext}
\usepackage{newpxmath}

% ── Colors ────────────────────────────────────────────────────────────────────
\usepackage[dvipsnames]{xcolor}
\definecolor{SectionBlue}{HTML}{1B3A5C}
\definecolor{AccentGray}{HTML}{4A4A4A}
\definecolor{QuoteBorder}{HTML}{8B9DAF}
\definecolor{RiskRed}{HTML}{C0392B}
\definecolor{RiskAmber}{HTML}{E67E22}
\definecolor{RiskGreen}{HTML}{27AE60}
\definecolor{BoxBg}{HTML}{F7F8FA}

% ── Tables ────────────────────────────────────────────────────────────────────
\usepackage{booktabs}
\usepackage{tabularx}

% ── Boxes ─────────────────────────────────────────────────────────────────────
\usepackage[framemethod=tikz]{mdframed}

% ── Lists ─────────────────────────────────────────────────────────────────────
\usepackage{enumitem}
\usepackage{needspace}

% ── Hyperlinks ────────────────────────────────────────────────────────────────
\usepackage{hyperref}
\usepackage{xurl}
\hypersetup{
  colorlinks=true,
  linkcolor=SectionBlue,
  urlcolor=SectionBlue,
  citecolor=SectionBlue,
  pdfauthor={Jim Freeman},
  pdftitle={crypt PR/FAQ},
  pdfsubject={Working Backwards: Product Discovery},
}

% ── Bibliography ─────────────────────────────────────────────────────────────
\usepackage[style=numeric,backend=biber,sorting=none]{biblatex}
\addbibresource{prfaq.bib}

% ── Section styling ───────────────────────────────────────────────────────────
\usepackage{titlesec}
\titleformat{\section}
  {\needspace{6\baselineskip}\Large\bfseries\color{SectionBlue}}{}{0pt}{}
\titleformat{\subsection}
  {\needspace{5\baselineskip}\large\bfseries\color{AccentGray}}{}{0pt}{}
\titleformat{\subsubsection}
  {\normalsize\bfseries\color{AccentGray}}{}{0pt}{}
\titlespacing*{\subsubsection}{0pt}{1em}{0.5em}[0pt]

% ── Document stage ───────────────────────────────────────────────────────────
\newcommand{\prfaqstage}[1]{\def\theprfaqstage{#1}}
\prfaqstage{hypothesis}

% ── Document version ─────────────────────────────────────────────────────────
\newcommand{\prfaqversion}[2]{\def\theprfaqmajor{#1}\def\theprfaqminor{#2}}
\prfaqversion{2}{3}

% ── Header/footer ────────────────────────────────────────────────────────────
\usepackage{fancyhdr}
\pagestyle{fancy}
\fancyhf{}
\fancyhead[L]{\footnotesize\textcolor{AccentGray}{Stage: \theprfaqstage{} \enspace\textbar\enspace v\theprfaqmajor.\theprfaqminor}}
\fancyhead[R]{\footnotesize\textcolor{AccentGray}{\thepage}}
\renewcommand{\headrulewidth}{0pt}

% ══════════════════════════════════════════════════════════════════════════════
\newcommand{\prsection}[1]{%
  \vspace{1em}{\large\bfseries\color{SectionBlue}#1}\vspace{0.3em}\par}

\newmdenv[topline=false,bottomline=false,rightline=false,linewidth=3pt,
  linecolor=QuoteBorder,backgroundcolor=BoxBg,innerleftmargin=12pt,
  innerrightmargin=12pt,innertopmargin=8pt,innerbottommargin=8pt,
  skipabove=\baselineskip,skipbelow=\baselineskip]{customerquote}

\newmdenv[topline=false,bottomline=false,rightline=false,linewidth=3pt,
  linecolor=SectionBlue,backgroundcolor=BoxBg,innerleftmargin=12pt,
  innerrightmargin=12pt,innertopmargin=8pt,innerbottommargin=8pt,
  skipabove=\baselineskip,skipbelow=\baselineskip]{spokespersonquote}

\newcounter{faqnum}
\newenvironment{faqpair}[1]{%
  \needspace{4\baselineskip}\vspace{0.5em}\refstepcounter{faqnum}%
  \noindent\textbf{\color{SectionBlue}Q\thefaqnum: #1}\par\nopagebreak[4]%
  \vspace{0.2em}\begin{adjustwidth}{1em}{0pt}\setlength{\parindent}{0pt}%
}{\end{adjustwidth}\vspace{0.5em}}
\newcommand{\faqref}[1]{\hyperref[#1]{FAQ~\ref*{#1}}}

\newcounter{featurenum}
\newcommand{\featureitem}[2]{%
  \refstepcounter{featurenum}%
  \item[\textbf{\color{SectionBlue}F\thefeaturenum.}]\textbf{#1}: #2}
\newcommand{\featureref}[1]{\hyperref[#1]{Feature~\ref*{#1}}}

\usepackage{changepage}

% ══════════════════════════════════════════════════════════════════════════════
\begin{document}

% ── Headline block ───────────────────────────────────────────────────────────
\begin{center}
  {\LARGE\bfseries\color{SectionBlue} Wizardry's Party of Six Returns, With Friends and AI Companions Who Remember You and Can Die} \\[0.8em]
  {\large\color{AccentGray} crypt fills every seat in the party with any mix of human players and companions, and with a model connected gives Zork and Wizardry fans a dungeon that never runs out} \\[0.3em]
  {\large\color{AccentGray} Built to work through Chip Huyen's AI-engineering curriculum end to end, with a game people enjoy as the proof}
\end{center}

\vspace{0.5em}

% ══════════════════════════════════════════════════════════════════════════════
\section*{Press Release}

SAN FRANCISCO: Today the crypt project released crypt, a text-adventure
dungeon crawler for players who grew up on Zork, Wizardry, and the old
MUDs. A party descends together, with up to six seats, filled in any mix
of friends and AI companions. Each companion has its own personality
and voice, remembers what happened in earlier sessions, and can die. A
deterministic engine enforces every rule for humans and companions
alike, the way a board game does. Reloading a save never brings a dead
companion back. Where the dungeon offers resurrection, a companion pays
the same price a human does. Under the party, with a model connected, a
language model writes a new dungeon for every descent, so the dungeon never
runs out (see \faqref{faq:party}).

\prsection{Problem}

The dungeon crawls that defined a genre ran out of content. Zork had to fit
a 32K microcomputer, so its authors cut it from 191 rooms to 110 and its
vocabulary from 911 words to 617 \cite{maher2012zilzmachine}. Wizardry
shipped ten fixed levels under one castle \cite{wikipedia2026wizardry1}.
Once a player solves them, the mystery is gone. The MUDs answered with crowd
authorship: LPMud wizards and LambdaMOO builders added rooms, monsters, and
objects from inside the running game
\cite{wikipedia2026lpmud,wikipedia2026lambdamoo}. Crowd authorship needs a
crowd, and a world grows only while its builders keep logging in. The party
had its own limit. Wizardry gave one player six characters to run alone
\cite{wikipedia2026wizardry1}. A MUD offered other people, but a party
still needed five of them free on the same evening. Games run entirely by a
language model promise endless content and then lose track of it. The top
models in RPGBench's evaluation write engaging stories while keeping the mechanics inconsistent
\cite{yu2025rpgbench}. AI Dungeon's own help page explains that older story
text is trimmed once the context fills \cite{latitude2025help}.

\prsection{Solution}

A player opens a party and fills its six seats in any mix. A player alone
starts with three companions and can add up to five. Two friends and three
companions make a full party, and so do six friends. A
companion is a character played by a language model, with a class, a
personality, and a voice of its own. It remembers the last descent: who took
the cursed ring, who ran from the ogre, who carried whom out of the pit. It
acts through the same moves a human player has, and the engine resolves its
sword swing or its spell with the same rules and the same dice. When a
companion dies, no reload undoes it; where the dungeon offers resurrection,
a companion pays the same price a human does. Under the party, with a model
connected, every descent after the first bundled one is a freshly generated
dungeon. Before a dungeon opens, the engine walks it to prove that
every locked door has a key and every goal can be reached. Around the text,
each companion speaks aloud in its own voice. Richer presentation follows as
each layer proves it can keep up without making the player wait: a score
that shifts with the scene, then panels that redraw around it
(see \faqref{faq:presentation}).

\prsection{Customer Quote}

\begin{customerquote}
\textit{``In Wizardry I ran all six characters myself. Now my son plays the
cleric, I play the fighter, and the other four are companions. Three
sessions ago I sold an amulet our thief, Wren, had found, to pay for my own
resurrection. Last night, splitting the loot from the lich's vault, she
brought it up and would not let it go. An hour later she held the bridge so
we could run, and the ogre killed her. The temple wanted more gold than the
party had, and reloading does nothing. She is gone. We have a new thief now,
and my son still won't talk to him.''}

\raggedleft \textbf{Dave Whitman}, lifelong text-adventure player
\end{customerquote}

\prsection{Getting Started}

A player installs crypt, runs it, and is in a bundled dungeon within a minute, with
three companions already in the party. Inviting a friend hands that
friend a seat. From inside Claude Code, the
game connects through a built-in bridge and Claude becomes the dungeon
master directly. From a terminal, the game runs its own client and talks to
whatever model the player has available, down to a fully offline mode that
still plays with built-in narration. A player joins the hosted service with
an invite link. There is no account at launch; an invite link is the
identity. There is no world to download. A dungeon is one command
away (see \faqref{faq:start}).

\prsection{Spokesperson Quote}

\begin{spokespersonquote}
``I built crypt to learn how to engineer an AI-enhanced product: where to
put the model, how to measure whether it is any good, and how to tune it. A
companion is the hardest exercise I could find: an agent with a memory and a
personality that has to play fair inside rules it cannot bend. The engine
knows which moves are legal and what happened, so I can score every
companion's rule-following against it, and every memory claim about
something the engine recorded. How well a
companion chooses, and whether it stays in character, I score with judges.
The game is the part people see; the learning is why it exists.''

\raggedleft \textbf{Jim Freeman}, crypt project
\end{spokespersonquote}

\prsection{Call to Action}

crypt is free and open source. A player can join a party on the hosted
crypt service with an invite link or run the game on their own machine. To play it, build it,
or read how the engine and the companions divide the work, visit the
repository at \mbox{github.com/punt-labs/crypt}.

\newpage

% ══════════════════════════════════════════════════════════════════════════════
\section*{Frequently Asked Questions}

\subsection*{External FAQs}

\begin{faqpair}{What is crypt and who is it for?}
crypt is a text-adventure dungeon crawler for players who loved Zork,
Wizardry, and the MUDs. Its center is the party: six seats, filled by
friends, by AI companions, or by both. Companions remember earlier sessions,
keep their own personalities, speak in their own voices, and can die, with
no reload to undo it (see \faqref{faq:death}). The classic mechanics run in a deterministic engine: grid dungeons,
turn-based combat, character classes, inventory, and leveling. With a model
connected, a language model writes a new dungeon for each descent, so the
party never runs out of places to go. In rules-only play the game serves
hand-authored scenarios from a bundled set (see \featureref{feat:party}).
\end{faqpair}

\begin{faqpair}{How is this different from AI Dungeon, or from other games that pair an engine with a language model?}
crypt plays at Wizardry's tempo: one keypress per decision, in a corridor
or in combat, with the party's state always on screen. Its terminal client
has no page to freeze, and vox plays each voice in a separate process. The
engine answers the moment a key is pressed, and narration arrives behind
it. The closest commercial neighbor, TableForge, plays as a browser chat
surface. In one free-tier session, a turn was a mouse path: open Quick
Actions, pick an action chip or type, press Send, and wait. By default voice
played only when the player clicked a play button under each message, party
state sat behind tabs, and nothing observed was keyboard-reachable. The wait
appeared to have two layers. One was model and voice synthesis time, the layer crypt's
predictive cache targets. The other was the browser laboring under the
page itself, and the tab froze once and had to be reloaded. The session was
slow enough to doubt it holds attention for hours
\cite{freeman2026tableforgehandson}.

AI Dungeon is collaborative storytelling with no rules engine. Its
multiplayer mode sets no limit on players in one story
\cite{latitude2026multiplayerfaq}, and the story lives in the model's
context, which trims older text as it fills \cite{latitude2025help}. In
crypt, hit points, inventory, and death live in the engine and do not drift
as a session grows long.

TableForge's makers say it gives ``final authority over dice, damage,
conditions, and resources to code.''
An AI dungeon master narrates on top, a party seats up to six human
players, and one to four AI companions join depending on the plan
\cite{tableforge2026home,tableforge2026faq,tableforge2026multiplayercompare}.
Friends \& Fables, Hidden Door, Auferet, and Questie each cover part of the
same ground \cite{fables2026home,hiddendoor2026home,auferet2026multiplayer,questie2026partymode}.
On the pages read in November 2026, none of six reviewed products stated
five things a crypt player gets. A crypt player can read the source and
play offline with no model at all. A player inside Claude Code brings their
own model, and it runs the game through the engine's tools. A companion who
says she was cheated out of an amulet three sessions ago is checked against
what the engine recorded. A companion who fell to the ogre stays dead
through the game's own save and load, and on the hosted service. A score and panels follow the scene as each passes
its gate. Hidden Door comes nearest on two of these. A world creator can set
characters to permadeath. Its Creator MCP lets Claude author worlds, and
Hidden Door's own story engine runs play \cite{hiddendoor2026home}.
At least eight open-source games also pair a rules engine with a language
model, among them kliuchnikovv/dnd, Mnehmos/rpg-mcp, open-dungeon, and a
fork of Lebbitheplow/open-game-master
\cite{kliuchnikovv2026dnd,mnehmos2026rpgmcp,slowbutfast2026opendungeon,alecvlad2026opengamemaster}
(see \faqref{faq:landscape}).
\end{faqpair}

\begin{faqpair}{Can I play with friends, or do I have to play alone?}
Both, in one party. A player alone starts with three companions and can add
up to five. A player and two friends leave three seats for companions. Six
friends fill all six and play with no companions at all. A friend can join
between descents and take a companion's seat, and the companion steps aside with its memory intact
(see \featureref{feat:multiplayer}).
\end{faqpair}

\begin{faqpair}{What happens when a companion dies?}\label{faq:death}
It stays dead, the same as a human player's character. The engine records
the death the moment it happens, and reloading a save cannot undo it. Where
the dungeon offers resurrection, at a temple or through a spell, it follows
the same rules and costs the same for humans and companions. A fallen
companion's seat opens for a new companion, a new character, or a friend.
The game's own save and load never undo a death, and the hosted service
enforces it. A local player with the source can copy a save directory, as
Wizardry players copied disks; the game does not try to stop that
(see \featureref{feat:companions}).
\end{faqpair}

\begin{faqpair}{Do I need an API key or an internet connection to play?}
No. crypt runs in tiers. A large model accepts the widest range of
free-text commands, writes the longest and most varied descriptions, and
gives companions the most detailed recall and dialogue. A small local model
still generates, narrates, and plays companions, with shorter descriptions
and a narrower range of commands it understands. With no model at all, the
game falls back to built-in rules and templates, companions fight by
scripted tactics, and the game stays fully playable offline.
\end{faqpair}

\begin{faqpair}{Will a companion or the narrator cheat, forget my inventory, or break the rules?}
No. A model never decides a game outcome. Every rule, every die roll, every
point of damage, and every saved game is handled by the engine in code. A
companion chooses what to try, the same as a human player, and the engine
decides what happens. The narrator describes the result, but it cannot
change your hit points, revive a dead companion, or hand out a sword the
dungeon does not hold (see \faqref{faq:line}).
\end{faqpair}

\begin{faqpair}{How do I get started?}\label{faq:start}
Three ways in. On the hosted service, a player opens an invite link and is
in a party with three companions. There is no account at launch; the invite
link is the identity. On a terminal, a player installs crypt from the
repository, runs it, and plays against whatever model is configured. That
can be a model provider key or a small local model. With no model at all,
the game plays fully offline with built-in narration and scripted companion
tactics. Inside
Claude Code, a player adds the crypt bridge and Claude becomes the dungeon
master, running the game through the engine's tools
(see \featureref{feat:bridge}). In every mode a bundled scenario opens
within a minute of starting. A generated dungeon opens after the engine's
solvability walk passes, and the validation run measures how long that
takes.
\end{faqpair}

\begin{faqpair}{How much does it cost?}\label{faq:cost}
Nothing. crypt is free and open source. Run locally, it uses the player's
own model provider key, a small local model, or no model at all. Any model
cost is the player's own, and it can be zero. The hosted service is free
and invite-only, and it runs under a monthly spending ceiling. At the
ceiling, the hosted game drops to the small-model tier, then to rules-only
play, and new invites pause until the next month. In rules-only play no
model generates dungeons, so the game serves hand-authored scenarios from
the bundled set, which is finite (see \featureref{feat:dist}).
\end{faqpair}

\begin{faqpair}{What happens to my data?}\label{faq:data}
What a player types and what each companion remembers are game state, and
the server stores them with the rest of the game. On the hosted service,
that text goes to the model provider the operator configures, because the
model reads it to interpret commands and play companions. In local mode it
goes wherever the player's own model runs, and with a local model or no
model it never leaves the player's machine. There are no accounts at
launch, so the hosted service collects no email address; an invite link is
the identity. The engine bounds what prompt injection can do to outcomes:
an injected instruction cannot reposition a companion, change a roll, or
alter inventory. What it can change is what a companion or the narrator
says. On the hosted service, generated text passes an output check before
it reaches other players, and the check is one of the Chapter 10
guardrails exercises \cite{huyen2025aie}. The hosted service is not designed for children, and anything
typed into it reaches a model. A family that wants everything kept on its
own machine can run crypt locally.
\end{faqpair}

\newpage
\subsection*{Internal FAQs}

\subsubsection*{Purpose \& Learning Objectives}

\begin{faqpair}{Why build this? What is the real objective?}\label{faq:why}
crypt is a build-to-learn project, roughly four parts learning to one part
game. The curriculum is Chip Huyen's \textit{AI Engineering}
\cite{huyen2025aie}. Chapters 3 through 10 each get a concrete exercise in
a running application instead of a toy notebook; Chapters 1 and 2 are
background reading (see \faqref{faq:curriculum}). Her central discipline is
evaluation-driven development \cite{huyen2025aie}. In crypt the
deterministic engine knows which moves are legal and what happened, so
rule-following and memory have a recorded answer to score against. The
dungeon crawl is the vehicle, and I grew up playing the games it draws
on; the understanding is the product. The finish line is a public, playable
game running in production on AWS, with per-call latency, tokens, cost,
eval score, and error rate recorded for every model call. Production work is
part of the curriculum: Huyen's chapters on inference optimization and on
monitoring and the feedback flywheel need a live system to practice on
(see \featureref{feat:aws}).
\end{faqpair}

\begin{faqpair}{Why is the party the center of the product?}\label{faq:party}
Three reasons: the player, the curriculum, and the evidence. Wizardry's
party of six is the part of the genre I miss most, and playing with friends
or with model-played companions is the hook I want for myself. In Huyen's
curriculum, a companion is the deepest agents exercise in the project: it
plans, calls tools, keeps a memory, and fails in ways that need guarding
against. Generative Agents is the peer-reviewed anchor for that design.
Park et al. gave 25 language-model agents memory, reflection, and planning,
and the agents behaved believably over days of simulated life
\cite{park2023generativeagents}. A candidate dataset exists. FIREBALL
holds about 25,000 sessions of D\&D 5e play recorded through a Discord bot,
with structured game state \cite{zhu2023fireball}. Whether it transfers to
how players act in a crypt party is untested, and its license terms for
training use are checked before any fine-tuning run.

Build order follows from the party. Companions come before panels, because
a beautifully presented room with nobody worth adventuring beside is the
failure to avoid. Generated scenarios supply the places the party goes
(see \featureref{feat:dm}). Presentation comes last and deepens only as far
as its gates allow (see \faqref{faq:presentation}).
\end{faqpair}

\begin{faqpair}{Where is the boundary between classical code and the language model?}\label{faq:line}
The engine owns anything that must be exact and repeatable: state
transitions, combat resolution, inventory limits, movement, death, and
save and load. The models own anything that benefits from creativity and
language: generating a scenario, describing a room, interpreting free-text
commands, and playing a companion. A companion acts only by calling the same
engine tools a human player's commands reach. It can choose to attack,
cast, flee, or argue over the loot, and the engine resolves the attempt by
the rules. A companion's memory and personality live outside the engine.
They decide what the companion tries; the engine decides what happens.
Generated presentation sits on the same side as narration: a panel, a
score, or a voice displays engine state and never changes it.

The evidence for keeping state in code is consistent. RPGBench finds that
the top models in its evaluation write engaging stories as game engines while leaving the
mechanics inconsistent \cite{yu2025rpgbench}. Song et al. show that giving
a game master function calls for state improves both narration and state
consistency \cite{song2024labyrinth}. Infobip reports taking dice and turn
order away from the model after its pure-language-model D\&D engine fell
below 65\% accuracy as context grew. That figure is a vendor self-report
with no published method \cite{infobip2025dnd}.

Other systems let a model enforce the rules. CriticaLLM pairs a game-master
model with a judge model that validates player input
\cite{berlaimont2026criticallm}. Bicking's Intra separates some state from
the model but leaves outcome resolution to the model
\cite{bicking2025intra}. A study of hybrid narrative managers at ICTAI 2026
compares several designs and finds that none dominates on every measure
\cite{clerc2026hybridnarrative}. crypt keeps rules in code because the
engine's answer can be checked. The same state and the same action give
the same result, and a test can assert it. A judge model's ruling has no
such check, so a wrong ruling looks the same as a right one. The engine
owning state is the
reason a companion's death stays final and its memory of an amulet matches
what the engine recorded. The exact placement of the line for each new
capability, such as how much a companion may propose before the engine
rejects it, is tested by building.
\end{faqpair}

\begin{faqpair}{How will the model's performance be evaluated?}\label{faq:eval}
With all three of Huyen's evaluation approaches \cite{huyen2025aie},
anchored where possible to what the engine records. Exact evaluation
covers three things. The interpreter is scored against a hand-labeled set
of player inputs and their intended actions. Agreement with the rule-based
parser is reported only on inputs the rule-based parser handles. Balance is
measured by running a fixed set of automated sessions, 200 per scenario,
against survival and progression targets. Companion rule-following is exact
too: every tool call a companion makes is checked against what the engine
allowed. Memory is exact for claims about recorded events, after one
extraction step. A judge model pulls factual claims out of companion
dialogue, such as who sold the amulet or who fled the ogre. Each extracted
claim is then checked exactly against the engine's event log. The exact
score inherits the extractor's recall, which is reported alongside it.

AI-as-a-judge covers what the engine cannot score: whether a companion's
choice was a good one, whether its personality holds across sessions, and
narration quality. Tester personas do the judging. Their scores are
treated cautiously, because Huyen shows that a judge's scores drift over
time \cite{huyen2025aie}. For generated panels, a judge also scores
multi-turn consistency: whether the screen still reads as the same game on
turn two hundred as on turn one. EvoGenUI-Bench measures that property for
generated interfaces \cite{evogenuibench2026}. Comparative evaluation ranks
one model or prompt against another, which Huyen finds more reliable than
absolute scores \cite{huyen2025aie}. These feed a private leaderboard that
ranks each model tier by cryptd's own criteria, so model and prompt choices
are measured before they ship (see \featureref{feat:eval}).
\end{faqpair}

\begin{faqpair}{What AI-engineering skills is this meant to build?}\label{faq:skills}
Seven, in order of depth. Integrating a model into a working application
behind a clean boundary. Designing fallback tiers so the game keeps every
rule and stays playable, with plainer narration, when the model is weak or
absent. Building evaluation that
scores model output against ground truth. Tuning prompts and swapping
models to raise those scores. Building agents with memory and personality
that act only through tools. Serving several models per turn within a
latency and cost budget, in production, with monitoring. Fine-tuning a
small model for the game's narrow narration, parsing, and companion tasks.
Each builds on the one before, and the game gives each a concrete place to
practice.
\end{faqpair}

\begin{faqpair}{How does the project map to Chip Huyen's curriculum?}\label{faq:curriculum}
Chapters 3 through 10 of \textit{AI Engineering} \cite{huyen2025aie} each
get a concrete exercise in crypt, and the milestones are sequenced to
match. Chapters 1 and 2, an introduction to building with foundation
models and an account of how those models work, are background reading
with no exercise of their own. Evaluation methodology
and evaluating AI systems (Ch 3 and 4) become the evaluation suite and the
private leaderboard that rank the model tiers. Prompt engineering (Ch 5)
becomes the interpreter and narrator system prompts, plus the persona and
defensive prompts that keep a companion in character and inside the rules,
scored for sensitivity when they change. Chapter 6 covers retrieval and
agents, and the companions exercise both. Companion memory is retrieval:
before a companion acts, it retrieves from the engine's recorded events
and from its own earlier reflections. Its agent side is planning and
action only through the engine's tools \cite{park2023generativeagents}.
The game-as-dungeon-master and the tester personas are agents over the
same tools. Fine-tuning (Ch 7) is an experiment, run whatever its outcome:
train one small model on the harvested corpus and compare it to prompting
on the leaderboard. A tuned model ships only if it wins that comparison.
Dataset engineering (Ch 8) becomes a corpus of intent-to-action,
event-to-narration, and companion-decision pairs, harvested from real and
automated sessions. FIREBALL, D\&D 5e play recorded through a Discord bot,
is a candidate addition \cite{zhu2023fireball}; its transfer to crypt is
untested, and its license terms for training use are checked before any
fine-tuning run. Inference
optimization (Ch 9) becomes the latency budget for a turn with several
companions and the predictive cache that prepares the next screen while
the player reads. Architecture and user feedback (Ch 10) becomes production
on AWS. Every model call records latency, tokens, cost, eval score, and
errors, and player comments and persona reports feed the next iteration.
\end{faqpair}

\subsubsection*{Audience \& Demand}

\begin{faqpair}{Who is the audience, and is demand a concern?}\label{faq:audience}
The audience is players nostalgic for Zork, Wizardry, and MUDs who miss
adventuring in a party and want a dungeon that does not run out. Demand is
a minor concern here by design. This is a learning project, roughly four
parts learning to one part game. Success is a public, playable game that
records per-call latency, tokens, cost, eval score, and error rate for
every model call. Demand still has a floor, because Huyen's chapter on
architecture and user feedback (Ch 10) needs players. Its monitoring and
feedback exercises need at least 10 returning players over a month. A
returning player is one saved party that plays on two or more distinct
days in a month; the save identity, not the invite link, is the unit.
Below that number, those chapters run on
persona play-testers, and the results are labeled as persona results
(see \faqref{faq:success}). Any wider demand claim is a
hypothesis with no primary research behind it yet. \texttt{[CITATION
NEEDED]}
\end{faqpair}

\subsubsection*{Competition}

\begin{faqpair}{Who else builds games like this, and where does crypt stand?}\label{faq:landscape}
Pairing a rules engine with a language model is an established pattern,
commercially and in open source.

Six commercial products were reviewed from their own pages in November
2026. TableForge is the closest neighbor. Its makers say it gives ``final
authority over dice, damage, conditions, and resources to code'' and uses
AI ``for interpretation and narration''
\cite{tableforge2026multiplayercompare}. A host's plan sets the party: up
to six human players and one to four ``travelling NPC companions that
fight beside you,'' with ``persistent cross-session memory'' claimed for
the campaign \cite{tableforge2026home,tableforge2026faq}. Friends \& Fables
seats up to six players around an AI game master that ``narrates,
adjudicates rules, and reacts to your world.'' Its home page names no AI
companions \cite{fables2026home}. Hidden Door runs a story
engine that checks every beat ``against canon,'' offers world creators a
``None / Normal / Permadeath'' setting, and has a Creator MCP for building
worlds from Claude \cite{hiddendoor2026home}. Auferet has a dice roller and
stat tracking, but ``the AI runs the world, the NPCs, and the outcomes.''
Its page warns that ``AI narration can still apply a rule incorrectly''
\cite{auferet2026multiplayer}. Questie's Party Mode seats one human with up
to four voiced AI companions and has no rules engine beyond a d20 roll
\cite{questie2026partymode}. AI Dungeon's multiplayer has no rules engine
and no set player limit \cite{latitude2026multiplayerfaq}.

One free-tier TableForge campaign was also played hands-on, in Safari on
macOS \cite{freeman2026tableforgehandson}. The engine owns the numbers, as
the vendor says: a 5e character sheet, a grid map, and a quest log sit in
tabs beside the chat. A turn is a run of clicks: Quick Actions, an action
chip or typed text, Send, then a play button under each message to hear
it. Group voice autoplay exists as a setting and defaults off. The session
felt sluggish, and the browser tab froze once and needed a reload.
Companions were not seated on the free tier played; the Companions tab
read ``No companions travel with the party yet.'' TableForge's companion
memory therefore remains a vendor claim, unobserved; no page addresses
companion death.

None of six reviewed products stated five properties of crypt. None stated
open-source code with a fully offline, rules-only mode; all six are hosted
browser or app services. None stated a bridge that lets the player's own
model, such as Claude Code over MCP, run the game. Hidden Door's Creator
MCP authors worlds, and its own engine runs play. None stated that
companion memory is checked against engine-recorded state; TableForge's
campaign memory is stored game state, and its pages do not say how
companions use it. None stated permanent death for companions that no
reload undoes; Hidden Door's permadeath setting applies to characters, and
its page does not address companions or reloads. None stated generated
panels or a layered score; TableForge and Friends \& Fables generate art,
battle maps, and voice, and Questie gives each companion its own voice.
Each absence rests on one to three vendor pages per product, so it records
what the vendors state.

In open source, at least eight projects keep dice, hit points, and state in
code and let a model narrate or propose actions
\cite{kliuchnikovv2026dnd,mnehmos2026rpgmcp,thegreystone2026rpgmcp,shadowgate152026dungeonmaster,slowbutfast2026opendungeon,alecvlad2026opengamemaster,datashaman2026aittrpg,undergroundrap2026aimud}.
kliuchnikovv/dnd is the nearest in form. It runs a Go engine in a terminal
that never calls a model, and it still plays offline with its narrator
disabled \cite{kliuchnikovv2026dnd}. Mnehmos/rpg-mcp runs as an MCP server
in Claude Code under the rule that models propose and never execute
\cite{mnehmos2026rpgmcp}. Their GitHub star counts are between 0 and 46
for the repositories whose counts were collected. The open-game-master
entry is a fork of Lebbitheplow/open-game-master
\cite{alecvlad2026opengamemaster}, so its star count is not meaningful. TaleWeave AI is a multiplayer text adventure with AI characters
whose model drives character actions with less engine authority
\cite{ssube2024taleweave}.

Outside games, four neighbors set expectations. Generative Agents
is the research reference for agents with memory, reflection, and planning,
set in a simulated town where nobody can die
\cite{park2023generativeagents}. NVIDIA ACE puts model-driven characters
into graphical games, including inZOI, where an on-device small language
model drives companion characters \cite{nvidia2025aceinzoi}. AI Dungeon
is a different category, collaborative storytelling with no rules engine,
with company-reported 8M+ registered players \cite{latitude2025growth}. Google's A2UI is the
architecture analogue for generated panels: a model emits a declarative
description from a trusted widget catalog, and the client renders it
\cite{google2025a2ui}. The survey by Gallotta et al. maps these roles,
player, character, game master, and designer, across the field
\cite{gallotta2024llmgames}.

TableForge, or any of the open-source projects, could add engine-checked
companion memory or permanent companion death, and crypt holds no
structural barrier against that. A
learning project does not need one. The party is where the curriculum goes
deepest, and that holds whatever other games ship.
\end{faqpair}

\subsubsection*{Technical}

\begin{faqpair}{What are the major technical risks?}\label{faq:risks}
Seven. Six are open risks on Must Do features, and each has a spike and a
gate before the work that depends on it.

Companion memory fidelity. A companion that misremembers who sold the
amulet breaks the promise. Memory is retrieval over the engine's recorded
events and the companion's own reflections, and every memory a companion
cites is checked against the event log. Spike: the validation run in
\faqref{faq:nextstep}. Gate: memory-claim accuracy of 95 percent; below 80
percent, retrieval is redesigned (see \faqref{faq:success}).

Per-turn latency and cost at three to five companions. Three companions run
three agents, plus the narrator, on every turn, and five run five.
Batching companion decisions into one call, running companions on a small
model, and letting them decide while the player reads are the planned
answers. Spike: measure p95 turn latency and cost per turn with three and
with five companions on each tier. Latency gate, for p95 with three
companions on the hosted tier: under 4 seconds is the target. Between 4 and
8 seconds, companion calls are batched into one request per turn before any
presentation tier ships. Above 8 seconds, the solo default drops to two
companions, batching is mandatory, and the panels tier stays blocked. Cost
gate: the same spike records tokens and cost per turn at three companions
on the hosted tier. That number times a monthly turn budget sets the hosted
spending ceiling, and the Viability rating moves off MEDIUM when the ceiling
is written into the deployment config.

Scenario solvability. A generated dungeon must have a key for every locked
door, a way past every trap, prices a party can pay, and cursed items whose
curse can be lifted. Spike: generate a batch of dungeons, run the engine's
solvability walk on each, and play a sample with automated parties. Gate:
every generated dungeon passes the walk before it opens, and a failure
falls back to a hand-authored scenario.

Multiplayer transport. The transport for a shared party is open (DES-034),
and the daemon's concurrent sessions are the starting point. Spike: two
clients in one party over the chosen transport, with turn order enforced by
the engine. Gate: no transport work starts until the single-companion
validation run passes.

Model quality for scenario generation on the small tier: whether a small or
local model can write a scenario and play a companion well enough, or
whether only a large model will do. The leaderboard answers this per tier.
Gate: a small-tier generated scenario passes the solvability walk on the
first try at least 70 percent of the time; below that, small-tier play uses
bundled scenarios.

The integration boundary: a companion will try moves the rules forbid, and
the engine has to refuse them without the refusal breaking character. Gate:
under 5 illegal attempts per 100 companion turns.

Generated presentation: latency and session-long acceptance of a screen
that changes every turn are both unanswered, and each presentation tier
waits on its own gate (see \faqref{faq:presentation}).
\end{faqpair}

\begin{faqpair}{How do companions run in each mode, and where do their model calls go?}\label{faq:modes}
Three modes, three answers. On the hosted service, the daemon runs every
companion on the model the operator configures. Every one of those calls
records latency, tokens, cost, eval score, and errors in the AWS
monitoring. In the local terminal client, the daemon runs companions on the
player's configured model or on the small local model. In the rules-only
tier, offline or hosted at the spending ceiling, companions follow scripted
tactics and make no model calls. No model generates dungeons there either,
so the game serves hand-authored scenarios from the bundled set, which is
finite. Through
the Claude Code bridge, Claude is the interpreter and the narrator, and the
companions run inside the same Claude session through the engine's tools.
Their model calls are the player's own, so the project's monitoring sees
only the engine events those calls produce and never the calls themselves.
Evaluation of bridge play uses the engine's event log alone.
\end{faqpair}

\begin{faqpair}{What are the dependencies and their fallbacks?}\label{faq:deps}
Six. vox, the Punt Labs text-to-speech tool, gives each companion a voice;
without it, the game is text only. ElevenLabs Music generates the layered
score; without it, there is no score. AWS is the deployment target for the
hosted service, and at hypothesis stage the project needs no second cloud
as a fallback. Model providers serve the large and small tiers; the tiered
design already falls back from a large model to a small local one to
rules-only play. Claude Code is one client among three, so the terminal
client and the hosted service work without it. The multiplayer transport
is open per DES-034, and the release ships with it. If the transport
slips, single-player with companions ships first, and friends join once
the transport works.
\end{faqpair}

\begin{faqpair}{How does generated presentation earn its place in the game?}\label{faq:presentation}
By passing a gate the project measures before the tier ships. Presentation
deepens in three tiers, in ascending risk.

Every tier is measured against tempo. TableForge's wait, from Send to text
and from Send to voice, is the baseline the predictive cache is measured
against; a follow-up timing run of ten turns will put a number on it
\cite{freeman2026tableforgehandson}. Click count per turn is the second
tempo metric: the target is one keypress per action and zero clicks to see
party state.

Voice comes first. Each companion speaks in its own voice through vox, the
text-to-speech tool Punt Labs already ships. TableForge voiced about four
long sentences per turn \cite{freeman2026tableforgehandson}. crypt's spoken
narration defaults to at most two sentences per turn, each under 15 words,
in the Zork register where a room is two lines and the prompt returns; a
``tell me more'' keypress requests the long version. Sentences per spoken
turn and words per sentence are the Chapter 5 prompt targets, synthesis
starts on the first sentence, and the gate is integration.

A layered adaptive score comes second: stems that enter and leave as the
scene shifts from corridor to combat to treasure room. Generation uses
ElevenLabs Music, with Audio Reference steering each new stem toward the
tempo, instrumentation, and mood of the ones already playing
\cite{elevenlabs2026music}. The vendor does not claim key matching, so
key is part of what the gate measures. Its gate is coherence: stems
generated across a full session still sound like one piece of music to a
listening judge.

Generated panels come third: the screen around the text redraws to fit the
scene, with a map, a party roster, or a ledger of loot. The model fills a
schema from a fixed widget vocabulary, the A2UI pattern, and the client
renders it \cite{google2025a2ui}. In pairwise tests, people prefer
generated pages to markdown answers \cite{leviathan2026generativeui}, and
they prefer generated task-specific interfaces in over 70\% of cases
\cite{chen2025generativeinterfaces}. Both studies measure one-shot pages
for information tasks, so their results do not transfer directly to a game
screen that persists across a session. Two risks remain open. Latency is the
documented failure: Gemini's dynamic view takes a minute or more to
generate \cite{google2025genui}. Session-long acceptance of a screen that
changes every turn has no measurement anywhere, and Nielsen Norman Group
names lost consistency as the cost of generated UI
\cite{nng2024generativeui}. The design answer to consistency is a fixed
frame: the layout holds still, and only the contents of its panels
regenerate. The answer to latency is a predictive cache. Turn-based play
gives the game dead time while the player reads, and the engine already
enumerates the legal next states, so the set of screens to prepare per turn
is small. The cost of a prepared panel is measured in the latency spike
for this tier, because no published figure covers schema-filled panels. The
gate for panels is a cache hit rate and a p95 latency under a budget fixed
before the tier is built.

Generated images in terminals that support them are a fourth, optional
step. Terminal image protocols are fragmented across Kitty, iTerm2, and
Sixel \cite{kitty2026graphics,iterm2images,arewesixelyet}, so the text and
half-block display stays the baseline, and the gate is reliable protocol
detection. For latency, Unbounded reached about one second per generated
scene image with distillation \cite{li2025unbounded}, which is the
reference point for this tier.
\end{faqpair}

\begin{faqpair}{What happens if generated presentation fails its gates?}\label{faq:fallback}
If two or more presentation tiers fail their gates, the project turns to an
agent-authored world. The part of crypt that sets it apart then has to come
from the dungeon growing, because presentation cannot supply it. AI wizards, in the LPMud
sense of builders with the power to add to the world
\cite{wikipedia2026lpmud}, dig new rooms, monsters, and puzzles overnight.
The engine validates every build for solvability and balance before it
opens, and companion parties play-test it before a human sees it. SpacetimeMUD
already has an agent building rooms through the player interface at demo
scale \cite{clockworklabs2024spacetimemud}. The fallback exercises the same
curriculum from another direction: planning agents, exact evaluation of
their output, and a feedback flywheel from player ratings. Its risk is
quality, as critiques of the agent-populated network Moltbook show
\cite{wikipedia2026moltbook}. It is the last Should Do item, and it starts
only on that trigger (see \featureref{feat:agentworld}).
\end{faqpair}

\begin{faqpair}{What is the state of the build, and what remains?}
Working today: the deterministic engine, the daemon that serves it with
concurrent sessions, the thin client with a terminal UI, the Claude Code
bridge, and the tiered model integration. vox, the voice tool companions
use, ships separately. The order of the remaining work is content first,
presentation second. The validation run comes first (see
\faqref{faq:nextstep}). Companions with memory and permanent death follow,
then generated scenarios with engine-proved solvability, then the hosted
service on AWS with monitoring, invite-only behind a monthly ceiling, and
the multiplayer transport, so friends join on launch day. The fine-tuning experiment runs once the
harvested corpus exists. Presentation follows in gate order: voice, music,
panels, images. If scope compresses, presentation tiers are cut from the
top down, and the party and production stay.
\end{faqpair}

\begin{faqpair}{What is the next step to validate the plan?}\label{faq:nextstep}
One companion, one scripted dungeon, repeated sessions, before any
multiplayer transport work. The run scores three things. Rule-following:
the engine accepts zero illegal moves, and the run counts how often the
companion attempts one. Memory: every claim the companion makes about an
earlier event is checked against the engine's event log. Latency: p95 turn
time on each model tier with one companion; the three-companion latency and
cost figures come from the spike in \faqref{faq:risks}. The thresholds in \faqref{faq:success} decide
whether companions go forward as designed, need re-prompting, or need a
retrieval redesign. This run is the first mission.
\end{faqpair}

\subsubsection*{Scope \& Success}

\begin{faqpair}{What does success look like for a learning project?}\label{faq:success}
Two things. A documented understanding of where the model helps, where it
falls short, how to measure its performance, and what it takes to improve
it. And a public, playable game running in production on AWS, recording
per-call latency, tokens, cost, eval score, and error rate, in which a
nostalgic player enjoys a party. The game is where those lessons get
tested. A working evaluation suite, a defensible answer to the
code-versus-model boundary, companions that keep their memories and their
rules, and hands-on experience with serving, monitoring, tuning, and
fine-tuning are the deliverables.

Five thresholds turn that into decisions.

Companion illegal moves. The engine accepts none, so the measure is
attempts per 100 companion turns. Target: under 5. Above that, the
companion prompts are reworked before companions go further.

Memory-claim accuracy against engine-recorded events. Target: 95 percent.
Below 80 percent, companion memory retrieval is redesigned.

p95 turn latency with three companions on the hosted tier. Target: under 4
seconds. Between 4 and 8 seconds, companion calls are batched into one
request per turn before any presentation tier ships. Above 8 seconds, the
solo default drops to two companions, batching is mandatory, and the
panels tier stays blocked.

Cost per turn. The latency spike (\faqref{faq:risks}) records tokens and
cost per turn at three companions on the hosted tier. That number times a monthly turn budget sets
the hosted ceiling, and the Viability rating moves off MEDIUM when the
ceiling is written into the deployment config.

Real players for the monitoring and feedback chapters. Minimum: 10
returning players over a month. A returning player is one saved party that
plays on two or more distinct days in a month; the save identity, not the
invite link, is the unit. Below that, those chapters run on persona
play-testers, and their results are labeled as persona results.
\end{faqpair}

% ══════════════════════════════════════════════════════════════════════════════
\newpage
\section*{Risk Assessment}

\renewcommand{\arraystretch}{1.6}
\begin{tabularx}{\textwidth}{@{} l l X @{}}
\toprule
\textbf{\color{SectionBlue}Risk} & \textbf{\color{SectionBlue}Rating} & \textbf{\color{SectionBlue}Assessment} \\
\midrule
Value      & \textcolor{RiskGreen}{LOW} & The project is four parts learning to one part game, so its value is the engineering knowledge gained, and the work produces that knowledge whether or not players arrive. Player demand has a floor: the monitoring and feedback chapters need 10 returning players over a month. Below that, they run on persona play-testers and report persona results (see \faqref{faq:audience}). \\
Usability  & \textcolor{RiskAmber}{MEDIUM} & Four open usability risks, each with a measurement. Parser frustration, the known failure of the genre: the interpreter tries the rules parser first and calls a model only for input the rules cannot parse. Measure: the share of player inputs the rules interpreter resolves without a model, target over 80 percent on the bundled scenarios. Companion chatter: one companion speaks unprompted per turn, and the persona play-testers rate readability. Waiting on agent calls: companions decide while the player reads, and the p95 latency gate in \faqref{faq:risks} sets what happens above 4 and 8 seconds. Prompt injection by other humans in the party: the engine bounds what injection can do to outcomes, so an injected instruction cannot reposition a companion, change a roll, or alter inventory. What it can change is what a companion or the narrator says. On the hosted service, generated text passes an output check before it reaches other players, one of the Chapter 10 guardrails exercises \cite{huyen2025aie}. \\
Feasibility & \textcolor{RiskAmber}{MEDIUM} & The engine, daemon, client, concurrent sessions, and model integration already work. Six Must Do risks are open, each with a spike and a gate: companion memory fidelity; per-turn latency and cost at three to five companions; scenario solvability for doors, keys, shops, traps, and cursed items; the multiplayer transport, open per DES-034, which the release ships with; model quality for small-tier scenario generation; and the integration boundary (see \faqref{faq:risks}). Presentation tiers each ship only through their own gate (see \faqref{faq:presentation}). \\
Viability  & \textcolor{RiskAmber}{MEDIUM} & No revenue model to satisfy, but the hosted service spends money on every model call. It launches invite-only with a monthly spending ceiling. At the ceiling, the service drops to the small-model tier, then to rules-only play, and new invites pause. The three-companion latency spike records tokens and cost per turn on the hosted tier; that number times a monthly turn budget sets the ceiling. The rating moves off MEDIUM when the ceiling is written into the deployment config (see \faqref{faq:cost} and \faqref{faq:success}). \\
\bottomrule
\end{tabularx}

% ══════════════════════════════════════════════════════════════════════════════
\newpage
\section*{Feature Appendix}

\subsection*{Must Do}
\begin{enumerate}[nosep,leftmargin=2.5em]
  \featureitem{Party of six}{any mix of human players and companions fills the six seats; a solo party starts with three companions and can seat up to five, and a friend can take a companion's seat between descents}\label{feat:party}
  \featureitem{Companions}{model-played characters with a personality and memory across sessions, acting only through the engine tools a human player uses; memory is retrieval over engine-recorded events and the companion's own reflections; death cannot be undone by reloading a save, and resurrection, where the dungeon offers it, follows the same rules and costs as for humans}\label{feat:companions}
  \featureitem{Multiplayer transport}{friends in different places share one party; the transport is open per DES-034, starting from the daemon's concurrent sessions, and ships in the release (see \faqref{faq:deps})}\label{feat:multiplayer}
  \featureitem{Deterministic engine}{rules, combat, inventory, movement, death, save and load, run in code and never in a model}\label{feat:engine}
  \featureitem{Daemon and tiered model integration}{large model, small local model, and a no-model fallback, so the game is always playable}\label{feat:tiers}
  \featureitem{Free-text interpretation}{the model turns plain-language input into engine actions}\label{feat:interpret}
  \featureitem{Generated narration}{the model describes rooms, events, and characters on demand}\label{feat:narrate}
  \featureitem{Claude Code bridge}{Claude plays the dungeon master through the stdio bridge into the engine, and companions run inside the same Claude session through the engine's tools (see \faqref{faq:modes})}\label{feat:bridge}
  \featureitem{Generated scenarios}{the model writes a new dungeon for each descent, with shops, traps, and cursed items, and the engine proves it solvable before it opens}\label{feat:dm}
  \featureitem{Evaluation suite and private leaderboard}{scored, repeatable measurement of interpreter accuracy against a hand-labeled intent set, companion rule-following and memory claims against engine-recorded events, judge-scored decision quality and personality consistency, and game balance}\label{feat:eval}
  \featureitem{Fine-tuning experiment}{train one small model on the harvested corpus and compare it to prompting on the leaderboard, whatever the outcome; this is Huyen's fine-tuning chapter in practice}\label{feat:ftexp}
  \featureitem{AWS production}{the game runs in production on AWS, recording per-call latency, tokens, cost, eval score, and error rate for every model call, the system Huyen's inference-optimization and monitoring chapters practice on}\label{feat:aws}
  \featureitem{Hosted service and public availability}{a stranger can join a hosted party with an invite link, with no account at launch, or install and play without a developer checkout; the hosted service launches invite-only under a monthly spending ceiling, and at the ceiling it drops to the small-model tier, then to rules-only play, and new invites pause}\label{feat:dist}
\end{enumerate}

\subsection*{Should Do}
\noindent{\small Presentation tiers ship in this order, lowest feasibility risk first, each when its gate passes (see \faqref{faq:presentation}). The agent-authored world comes last, on its trigger.}
\begin{enumerate}[nosep,leftmargin=2.5em]
  \featureitem{Per-companion voice}{each companion speaks in its own voice through vox; gate: integration only}\label{feat:voice}
  \featureitem{Layered adaptive music}{stems that follow the scene, generated with ElevenLabs Music under Audio Reference; gate: stems stay coherent across a full session}\label{feat:music}
  \featureitem{Fixed-frame generated panels}{the layout holds still and panel contents regenerate from a fixed widget vocabulary, prepared ahead by a predictive cache; gate: cache hit rate and p95 latency under a stated budget}\label{feat:panels}
  \featureitem{Generated images}{scene images in terminals that support an image protocol, with text and half-block display as the baseline; gate: reliable protocol detection across fragmented terminal support}\label{feat:images}
  \featureitem{Fine-tuned small model shipped}{a model tuned for the game's narrow narration, parsing, and companion tasks ships only if the fine-tuning experiment shows it beats prompting and a stock small model on the leaderboard (see \featureref{feat:ftexp})}\label{feat:finetune}
  \featureitem{Agent-authored persistent world}{AI wizards that dig new rooms overnight, validated by the engine and play-tested by companion parties; pursued if two or more presentation tiers fail their gates, because the dungeon's growth then has to set crypt apart (see \faqref{faq:fallback})}\label{feat:agentworld}
\end{enumerate}

\subsection*{Won't Do}
\begin{enumerate}[nosep,leftmargin=2.5em]
  \featureitem{Model-decided game rules}{no model, narrator or companion, ever resolves combat, damage, or death; that is the one boundary the design refuses to cross}\label{feat:norules}
  \featureitem{Graphical 3D client}{the game is text-first by identity; a rendered 3D world is a different product}\label{feat:no3d}
  \featureitem{Real-time action}{crypt is turn-based like its ancestors; twitch gameplay is out of scope}\label{feat:norealtime}
  \featureitem{Monetization}{a build-to-learn project has no store, subscription, or paid tier}\label{feat:nomoney}
\end{enumerate}

% ══════════════════════════════════════════════════════════════════════════════
\newpage
\printbibliography

\end{document}
